\chapter{Complexity — Parse Cmp}
%%% generated by the blueprint pipeline from TCSlib.Complexity.SpaceComplexity.Machines.ParseCmp
\section{Overview}

Declarations of the Lean module \texttt{TCSlib.Complexity.SpaceComplexity.Machines.ParseCmp}.

\section{Declarations}

\begin{definition}[Action for the first prover skip]
\lean{Complexity.LogProg.pskip1Act}
\label{Complexity.LogProg.pskip1Act}
\leanok
\uses{Complexity.LogProg.xAct, Turing.Action}
Given a register index $r$, states $j_F,j_T$, and an optional Boolean answer, define an
action that uses the action prescribed by $r$ and $j_T$ when the answer is true, and the
action prescribed by $r$ and $j_F$ otherwise, including when the answer is absent.
\end{definition}

\begin{definition}[Action for the second prover skip]
\lean{Complexity.LogProg.pskip2Act}
\label{Complexity.LogProg.pskip2Act}
\leanok
\uses{Complexity.LogProg.xAct, Turing.Action}
Given a register index $r$, states $j_1,j_C$, a Boolean $b_1$, and an optional Boolean
answer, define an action that uses the action prescribed by $r$ and $j_C$ when $b_1$ is
false and the answer is true; in every other case, use the action prescribed by $r$ and
$j_1$.
\end{definition}

\begin{lemma}[Input symbol in a paired encoding]
\lean{Complexity.LogProg.inSym_pairInput}
\label{Complexity.LogProg.inSym_pairInput}
\leanok
\uses{Complexity.LogProg.dbl_replicate, Complexity.LogProg.inSym, Complexity.LogProg.inSym_succ, Turing.dbl, Turing.pairEncode, Turing.pairEncode_eq_dbl}
\difficulty{3}
For all natural numbers $n,k$ and Boolean words $u,w$, the symbol at one-based input
position $k+1$ in the encoding of the pair whose first component is $n$ copies of $1$
and whose second component encodes $(u,w)$ equals the symbol at zero-based position $k$
in the concatenation of $2n$ copies of $1$, the separator $01$, the doubled word $u$,
another separator $01$, and $w$.
\end{lemma}

\begin{lemma}[Paired-word skip execution]
\lean{Complexity.LogProg.pairSkip_run}
\label{Complexity.LogProg.pairSkip_run}
\leanok
\uses{Complexity.LogProg.Reaches, Complexity.LogProg.Reaches.refl, Complexity.LogProg.Reaches.step, Complexity.LogProg.inSym, Complexity.LogProg.pskip2Act, Complexity.LogProg.reaches_right, Complexity.LogProg.rrun, Complexity.LogProg.xCfg, Turing.Cfg, Turing.dbl}
\difficulty{5}
Let $P$ be a register program with oracle, let $r$ be a register, and let
$j_{P1},j_{P2F},j_{P2T},j_C$ be states. Assume that at every input and work-tape
reading, the transition at $j_{P1}$ is the prescribed first skip action on register $r$,
and the transitions at $j_{P2F}$ and $j_{P2T}$ are the prescribed second skip actions on
register $r$, with the false and true cases respectively; assume also that none of these
three states is an oracle-call state. For any configuration $c$, integer $p$, Boolean
word $u$, and natural number $q$ satisfying $q+2|u|+2\le |x|+1$, suppose that for every
$k<2|u|+2$ the input symbol at position $q+k$ equals the $k$th symbol of the doubled
word $u$ followed by $01$ (with out-of-range symbols interpreted as absent). Starting
from the configuration obtained from $c$ by setting the state to $j_{P1}$, the
input-head position to $q$, and register $r$'s head position to $p$, the program reaches
the configuration obtained from $c$ by setting the state to $j_C$ and the input-head
position to $q+2|u|+2$, with register $r$'s head position still $p$. At every earlier
time in this run, the configuration is obtained from $c$ by changing only the state and
input-head position and setting register $r$'s head position to $p$, and that state is
not an oracle-call state.
\end{lemma}

\begin{definition}[Action for the first decider]
\lean{Complexity.LogProg.dcmp1Act}
\label{Complexity.LogProg.dcmp1Act}
\leanok
\uses{Complexity.LogProg.xAct, Turing.Action}
Given a register index $r$, states $j_F,j_T$, and an optional Boolean answer, define an
action that uses the action prescribed by $r$ and $j_T$ when the answer is true, and the
action prescribed by $r$ and $j_F$ otherwise, including when the answer is absent.
\end{definition}

\begin{definition}[Action for the second decider]
\lean{Complexity.LogProg.dcmp2Act}
\label{Complexity.LogProg.dcmp2Act}
\leanok
\uses{Complexity.LogProg.xAct, Turing.Action}
Given a register index $r$, a state $j_1$, a state-valued function $j_{RB}$ on Booleans,
a Boolean $b_1$, and optional Boolean values $a,c$, define an action as follows. If
$b_1$ is false and $a$ is true, use the action prescribed by $r$ and $j_{RB}$ evaluated
at whether $c$ is absent; otherwise, if $c=b_1$, use the action prescribed by $r$ and
$j_1$; in all other cases use the action prescribed by $r$ and $j_{RB}$ evaluated at
false.
\end{definition}

\begin{lemma}[Doubled-word comparison run]
\lean{Complexity.LogProg.cmpDbl_run}
\label{Complexity.LogProg.cmpDbl_run}
\leanok
\uses{Complexity.LogProg.apply_xAct, Complexity.LogProg.dcmp2Act, Complexity.LogProg.inSym, Complexity.LogProg.moveInputPos_pos_val, Complexity.LogProg.reaches_right, Complexity.LogProg.rrun, Complexity.LogProg.rrun_add, Complexity.LogProg.rrun_one_x, Complexity.LogProg.xCfg, Complexity.LogProg.xCfg_read, Turing.Cfg, Turing.FinTM.bufferTape, Turing.dbl}
\difficulty{6}
Let $P$ be a register program with oracle, let $r$ be a register, and let
$j_{D1},j_{D2F},j_{D2T}$ and $j_{RB}(b)$ be states. Suppose the ordinary transitions at
these states are respectively the specified first comparison action, the second
comparison action for a false or true comparison bit, and the corresponding return
action; suppose none of these states is a call state. Let $c$ be a configuration whose
register $r$ stores the word $w_r$ from cell $0$, and let $w_i,q_0$ be a bit word and a
natural number such that $q_0+2|w_i|+2\le |x|+1$ and the input symbols at positions
$q_0+k$ agree, for every $k<2|w_i|+2$, with the $k$th symbol of the doubled word
$\operatorname{dbl}(w_i)$ followed by $01$. For every decomposition $w_i=p u_1$, $w_r=p
u_2$ and input-head position $ip$ with $ip=q_0+|p|$, there are natural numbers $T,k$ and
an input-head position $ip_k$ such that $|p|\le k\le |w_r|$, and the run from state
$j_{D1}$, head position $ip$, and register head at $|p|$ ends after $T$ steps in state
$j_{RB}(\mathbf 1_{w_i=w_r})$, at input position $ip_k$ and register-head position
$k-1$. At each earlier time $t<T$, it is in some state $s$ that is not a call state,
with register-head position $q$ satisfying $0\le q\le |w_r|$.
\end{lemma}

\begin{lemma}[Run across a paired-input prefix]
\lean{Complexity.LogProg.skipPrefix_run}
\label{Complexity.LogProg.skipPrefix_run}
\leanok
\uses{Complexity.LogProg.RProg, Complexity.LogProg.Reaches, Complexity.LogProg.Reaches.refl, Complexity.LogProg.Reaches.step, Complexity.LogProg.Reaches.trans, Complexity.LogProg.inSym, Complexity.LogProg.inSym_succ, Complexity.LogProg.reaches_right, Complexity.LogProg.rrun, Complexity.LogProg.scanR, Complexity.LogProg.skipAct, Complexity.LogProg.xAct, Complexity.LogProg.xCfg, Turing.Cfg}
\difficulty{4}
Let $P$ be a register program with oracle, let $r$ be a register, and let
$j_K,j_{K2},n_x$ be states. Suppose that at $j_K$ the transition scans the input and
advances the head on register $r$ by one, moving to $j_K$ when the scanned bit is true
and to $j_{K2}$ otherwise; at $j_{K2}$ it advances that register head by one and moves
to $n_x$. Suppose neither $j_K$ nor $j_{K2}$ is an oracle-call state. If the input is
$1^{2n}01$ followed by a word $\mathrm{rest}$, then the run from the configuration with
the tapes and output of $c$, state $j_K$, input-head position $1$, and register-head
position $p$ reaches the configuration with the same tapes and output, state $n_x$,
input-head position $3+2n$, and register-head position $p$; before that endpoint, each
run state has the form specified by the reachability relation, with no oracle call.
\end{lemma}

\begin{definition}[Bounded register reachability]
\lean{Complexity.LogProg.ReachesB}
\label{Complexity.LogProg.ReachesB}
\leanok
\uses{Complexity.LogProg.RProg, Complexity.LogProg.rrun, Complexity.LogProg.xCfg, Turing.Cfg}
Let $P$ be a register program with oracle $	ext{oracle}$, and let $c_0,c_1,c$ be
configurations on input $x$, with register $r$ and integer bound $L$. The predicate
holds if there is a time $T$ such that running $P$ with the oracle from $c_0$ for $T$
steps gives $c_1$, and for every $t<T$ there are a state $s$, input-head position $ip$,
and integer $q$ such that the configuration at time $t$ equals the configuration
obtained by replacing the corresponding register data of $c$ by $s,ip,r,q$, the program
has no call at $s$, and $-1\le q\le L$.
\end{definition}

\begin{lemma}[Bounded-position reaching]
\lean{Complexity.LogProg.Reaches.toB}
\label{Complexity.LogProg.Reaches.toB}
\leanok
\uses{Complexity.LogProg.RProg, Complexity.LogProg.Reaches, Complexity.LogProg.ReachesB, Turing.Cfg}
\difficulty{3}
Let $P$ be a register program with oracle, and let $c_0,c_1,c$ be configurations on the
same input, with register $r$ and integers $p,L$. If $c_0$ reaches $c_1$ relative to
$c,r,p$ and $-1\le p\le L$, then throughout that reaching computation, each intermediate
configuration has the head of register $r$ at some position $q$ with $-1\le q\le L$, and
the computation still reaches $c_1$.
\end{lemma}

\begin{lemma}[Transitivity of bounded reaching]
\lean{Complexity.LogProg.ReachesB.trans}
\label{Complexity.LogProg.ReachesB.trans}
\leanok
\uses{Complexity.LogProg.RProg, Complexity.LogProg.ReachesB, Complexity.LogProg.rrun_add, Turing.Cfg}
\difficulty{4}
Let $P$ be a register program with oracle, and let $c_0,c_1,c_2,c$ be configurations on
the same input, with register $r$ and integer bound $L$. If $c_0$ reaches $c_1$ and
$c_1$ reaches $c_2$, each relative to $c,r,L$ with the register head staying between
$-1$ and $L$ at every intermediate time, then $c_0$ reaches $c_2$ with the same bound
and relative to the same reference configuration and register.
\end{lemma}

\begin{lemma}[Indexing an Appended List]
\lean{Complexity.LogProg.getElem?_append_len}
\leanok
\difficulty{3}
For finite Boolean words $pre$ and $rest$ and any natural number $k$, the possibly
undefined entry of their concatenation at index $|pre|+k$ equals the possibly undefined
entry of $rest$ at index $k$.
\end{lemma}

\begin{lemma}[Run of the second-payload equality test]
\lean{Complexity.LogProg.jeqPairSnd_run}
\label{Complexity.LogProg.jeqPairSnd_run}
\leanok
\uses{Complexity.LogProg.Reaches.toB, Complexity.LogProg.ReachesB, Complexity.LogProg.ReachesB.trans, Complexity.LogProg.cmpPlain_run, Complexity.LogProg.cmpReturn, Complexity.LogProg.dbl_replicate, Complexity.LogProg.getElem?_append_len, Complexity.LogProg.inSym, Complexity.LogProg.inSym_succ, Complexity.LogProg.pairSkip_run, Complexity.LogProg.pskip1Act, Complexity.LogProg.pskip2Act, Complexity.LogProg.rrun_add, Complexity.LogProg.skipPrefix_run, Complexity.LogProg.xAct, Complexity.LogProg.xCfg, Turing.Cfg, Turing.FinTM.bufferTape, Turing.dbl, Turing.pairEncode, Turing.pairEncode_eq_dbl}
\difficulty{6}
Let a register program and oracle be fixed, with register $r$ and control states as in
the hypotheses. Suppose the specified transitions implement the second-payload
comparison routine, and the stated control states are not oracle-call states. If the
input is the self-delimiting encoding of a unary word $1^n$ followed by the encoding of
two bit strings $u,w$, and register $r$ stores the canonical binary expansion of a
natural number $a$, then the program, started in the comparison routine's initial state
with input-head position $1$ and register head at $0$, reaches a configuration in the
affirmative state exactly when $w$ is that binary expansion, and otherwise in the
negative state. The register head is at $0$ in the final configuration, and at every
earlier step the routine remains in a non-call state with that head between $-1$ and the
length of the binary expansion.
\end{lemma}

\begin{lemma}[Run of the first-payload equality test]
\lean{Complexity.LogProg.jeqPairFst_run}
\label{Complexity.LogProg.jeqPairFst_run}
\leanok
\uses{Complexity.LogProg.Reaches.toB, Complexity.LogProg.ReachesB, Complexity.LogProg.ReachesB.trans, Complexity.LogProg.cmpDbl_run, Complexity.LogProg.cmpReturn, Complexity.LogProg.dbl_replicate, Complexity.LogProg.dcmp1Act, Complexity.LogProg.dcmp2Act, Complexity.LogProg.getElem?_append_len, Complexity.LogProg.inSym_succ, Complexity.LogProg.rrun_add, Complexity.LogProg.skipPrefix_run, Complexity.LogProg.xAct, Complexity.LogProg.xCfg, Turing.Cfg, Turing.FinTM.bufferTape, Turing.dbl, Turing.pairEncode, Turing.pairEncode_eq_dbl}
\difficulty{6}
Let a register program and oracle be fixed, with register $r$ and control states as in
the hypotheses. Suppose the specified transitions implement the first-payload comparison
routine, and the stated control states are not oracle-call states. If the input is the
self-delimiting encoding of a unary word $1^n$ followed by the encoding of two bit
strings $u,w$, and register $r$ stores the canonical binary expansion of a natural
number $a$, then the program, started in the comparison routine's initial state with
input-head position $1$ and register head at $0$, reaches a configuration in the
affirmative state exactly when $u$ is that binary expansion, and otherwise in the
negative state. The register head is at $0$ in the final configuration, and at every
earlier step the routine remains in a non-call state with that head between $-1$ and the
length of the binary expansion.
\end{lemma}
